\section{Discussion}

LETT-NeXt demonstrates that sparse RECIST markers can support efficient 3D lesion segmentation under CPU inference constraints. 
The model generally localized the marked lesions, while the remaining errors mainly involved lesion extent and boundary placement. 
This pattern is consistent with the lower NSD than DSC on the public validation set and with the under-segmentation examples in \Cref{fig:percentile_examples}. 
In practice, this suggests that the RECIST marker provides strong localization information, but does not fully resolve the 3D lesion boundary.

The controlled ablations suggest that backbone choice was the main driver of performance. 
The MedNeXt-v2 f32 backbone provided the clearest performance gain, whereas auxiliary supervision acted more like weak training-time regularization than a definitive improvement. 
AutoZoom was neutral on average, but remained useful as a selective fallback for cases where the initial regional crop may not fully contain the lesion.

The main limitations are the small 49-case validation set, threshold selection on the same local validation split used for model development, limited hidden-test detail, and the absence of external validation. 
The ablations also tested components mostly in isolation, so interactions between backbone choice, auxiliary supervision, crop strategy, and postprocessing remain uncertain. 
Future work should evaluate LETT-NeXt on external datasets and improve boundary refinement. 
A promising direction is to pre-train a larger anatomy-aware model on dense multi-class anatomy--tumor labels, for example, from PanTS, before fine-tuning it for RECIST-marker tumor segmentation or distilling it into a compact student model. 
This strategy could improve anatomical context and boundary modeling while preserving efficient CPU inference.

Overall, LETT-NeXt shows that RECIST markers can support accurate and efficient local 3D lesion segmentation, but further work is needed to improve boundary accuracy and robustness across datasets.

\section*{Acknowledgements}
The authors express their appreciation to Novartis Norge AS and Akershus University Hospital for funding this work. The submitted method for the \emph{Foundation Models for Pan-cancer Segmentation in CT Images} competition is fully automatic and does not require manual intervention. Only data permitted by the organizers was used. We thank the challenge organizers, data providers, and CodaBench~\cite{xuCodabenchFlexibleEasytoUse2022} for hosting the competition.


